% GENERATED FILE. Edit canonical lessons, then run npm run generate:books.
% canonical-source-hash: fnv1a64:1ba9a95e69920ef8
% canonical-lessons: UR-C272-dar, UR-W272-ddal, UR-C272-main-khush-hun, UR-C272-banna, UR-C272-ummid, UR-C272-banna-chahti-hun, UR-R272-feelings-and-hopes

\chapter{Feelings and Hopes}
\label{ch:ur-feelings-and-hopes}
\hlchaptermodality{272}

\begin{chapteropening}
\textbf{By the end of this chapter:} \emph{I can say how I feel or felt (maiṅ khush hūṅ, mujhe ḍar lagā), what I want to become (maiṅ ustād bannā chāhtī hūṅ) and what I hope (mujhe ummīd hai ki …).}

Complete UR-R272-feelings-and-hopes: say how you feel, what you want to become, and what you hope.
\end{chapteropening}

\section[\texorpdfstring{\ur{ڈر} (ḍar)}{ڈر (ḍar)}]{\ur{ڈر} (ḍar) — fear}

\label{lesson:UR-C272-dar}



\begin{quote}
Before the new one: say \textquotedblleft{}I'm hungry\textquotedblright{}. (\emph{mujhe bhūk lagī hai}.) Fear is said the same way.
\end{quote}

\subsection*{You'll want to know: \ur{ڈر}}
\textbf{\ur{ڈر}} — \emph{ḍar} — \textquotedblleft{}fear\textquotedblright{}. It is masculine.

\begin{quote}
\textbf{\ur{مجھے} \ur{ڈر} \ur{لگتا} \ur{ہے۔}}
\end{quote}

\begin{quote}
\emph{mujhe ḍar lagtā hai.} — \textquotedblleft{}I'm afraid.\textquotedblright{}
\end{quote}

Fear strikes you, as hunger does: \textbf{\ur{مجھے}}, and \textbf{\ur{لگتا}} (\emph{lagtā}) for a masculine \textbf{\ur{ڈر}}. What you fear takes \textbf{\ur{سے}} (\emph{se}):

\begin{quote}
\textbf{\ur{مجھے} \ur{کتوں} \ur{سے} \ur{ڈر} \ur{لگتا} \ur{ہے۔}}
\end{quote}

\begin{quote}
\emph{mujhe kuttoṅ se ḍar lagtā hai.} — \textquotedblleft{}I'm afraid of dogs.\textquotedblright{}
\end{quote}

\begin{cousinweb}
\textbf{\ur{ڈر}} (\emph{ḍar}) is inherited. CDIAL traces it to Sanskrit \emph{dara}, \textquotedblleft{}fear\textquotedblright{} (CDIAL 6186), and lists the Hindi form \emph{ḍar}, with the curled-back \emph{ḍ}. Its first letter is new to your hand; the next lesson takes it out.
\end{cousinweb}

\subsection*{Your turn}
Say these aloud:

\begin{itemize}
\raggedright
  \item \emph{ḍar}, then \emph{mujhe ḍar lagtā hai}
  \item \textquotedblleft{}She got scared\textquotedblright{} — \emph{use ḍar lagā}
\end{itemize}

\begin{tcolorbox}[breakable,colback=teal!4,colframe=teal!35!black,title={Before you move on}]
What does \textbf{\ur{ڈر}} mean? (\textquotedblleft{}Fear\textquotedblright{}.) \textquotedblleft{}I'm afraid of dogs\textquotedblright{}? (\textbf{\ur{مجھے} \ur{کتوں} \ur{سے} \ur{ڈر} \ur{لگتا} \ur{ہے}}.)
\end{tcolorbox}
\section[\texorpdfstring{\ur{ڈ} (ḍāl)}{ڈ (ḍāl)}]{\ur{ڈ} (ḍāl) — the letter that opens \ur{ڈر}}

\label{lesson:UR-W272-ddal}



\begin{quote}
Before the new one: write \ur{ث} once, and name it. (\emph{se}.)

Now say \textbf{\ur{ڈر}} (\emph{ḍar}), \textquotedblleft{}fear\textquotedblright{}. This lesson takes one letter out of it and puts it in your hand.
\end{quote}

\subsection*{Script you'll notice: \ur{ڈ}}
\textbf{\ur{ڈ}} — \emph{ḍāl}, the sound \emph{ḍ}, made with the tongue curled back, as for \textbf{\ur{ٹ}} (\emph{ṭe}) and \textbf{\ur{ڑ}} (\emph{ṛe}).

It opens \textbf{\ur{ڈر}} (\emph{ḍar}). You already write \textbf{\ur{د}}: \textbf{\ur{ڈ}} is that shape wearing the small mark that \textbf{\ur{ٹ}} and \textbf{\ur{ڑ}} wear above them. Like \textbf{\ur{د}}, it does not join the letter after it.

\subsection*{Writing: \ur{ڈ}}
Put your pen on \ur{ڈ} and follow its line. Copy the shape you can see — slowly, and larger than it is printed.

\begin{quote}
This book does not yet tell you \textbf{where to start this letter or which way to travel}. It is not written down here until it can be written down with a source. Copying what is in front of you needs no such source, and it is how the shape gets into your hand in the meantime.
\end{quote}

\subsection*{Your turn}
\begin{itemize}
\raggedright
  \item \emph{Look:} at \textbf{\ur{ڈر}} and point at \ur{ڈ} inside it
  \item \emph{Trace it:} \ur{ڈ} three times, saying \emph{ḍ} as you finish each one
  \item \emph{Write it:} \textbf{\ur{ڈر}} once, with the small mark above its first letter
  \item \emph{Read it:} \textbf{\ur{ڈر}} aloud once more
\end{itemize}

\begin{tcolorbox}[breakable,colback=teal!4,colframe=teal!35!black,title={Before you move on}]
Which letter is this — \ur{ڈ}? (\textbf{ḍāl}, \emph{ḍ}.) What tells it from \ur{د}? (The small mark that \ur{ٹ} wears.) Which word did it come from? (\textbf{\ur{ڈر}}, \emph{ḍar}, \textquotedblleft{}fear\textquotedblright{}.)
\end{tcolorbox}
\section[\texorpdfstring{\ur{میں} \ur{خوش} \ur{ہوں} (maiṅ khush hūṅ)}{میں خوش ہوں (maiṅ khush hūṅ)}]{\ur{میں} \ur{خوش} \ur{ہوں} (maiṅ khush hūṅ) — how I feel}

\label{lesson:UR-C272-main-khush-hun}



\begin{quote}
Say \textquotedblleft{}I'm afraid\textquotedblright{}. (\emph{mujhe ḍar lagtā hai}.) Then \textquotedblleft{}happy\textquotedblright{} and \textquotedblleft{}sad\textquotedblright{}. (\emph{khush}, \emph{udās}.) Urdu says a feeling two ways.
\end{quote}

\subsection*{You'll want to know: I am happy, she was sad}
\begin{quote}
\textbf{\ur{میں} \ur{خوش} \ur{ہوں۔}}
\end{quote}

\begin{quote}
\emph{maiṅ khush hūṅ.} — \textquotedblleft{}I'm happy.\textquotedblright{}
\end{quote}

\begin{quote}
\textbf{\ur{کل} \ur{وہ} \ur{اداس} \ur{تھی۔}}
\end{quote}

\begin{quote}
\emph{kal voh udās thī.} — \textquotedblleft{}Yesterday she was sad.\textquotedblright{}
\end{quote}

\begin{quote}
\textbf{\ur{میں} \ur{پریشان} \ur{تھا۔}}
\end{quote}

\begin{quote}
\emph{maiṅ pareshān thā.} — \textquotedblleft{}I was worried.\textquotedblright{}
\end{quote}

The feeling word stays as it is; \textbf{\ur{ہوں}}, \textbf{\ur{تھی}}, \textbf{\ur{تھا}} follow the person. With \textbf{\ur{ہوا}} (\emph{huā}) the feeling arrives:

\begin{quote}
\textbf{\ur{یہ} \ur{سن} \ur{کر} \ur{وہ} \ur{خوش} \ur{ہوئی۔}}
\end{quote}

\begin{quote}
\emph{yih sun kar voh khush huī.} — \textquotedblleft{}Hearing this, she was pleased.\textquotedblright{}
\end{quote}

\subsection*{You'll want to know: what strikes you}
\begin{quote}
\textbf{\ur{مجھے} \ur{ڈر} \ur{لگا۔}}
\end{quote}

\begin{quote}
\emph{mujhe ḍar lagā.} — \textquotedblleft{}I got scared.\textquotedblright{}
\end{quote}

\begin{quote}
\textbf{\ur{مجھے} \ur{یہ} \ur{شہر} \ur{اچھا} \ur{لگا۔}}
\end{quote}

\begin{quote}
\emph{mujhe yih shahr achchhā lagā.} — \textquotedblleft{}I liked this city.\textquotedblright{}
\end{quote}

A feeling that strikes you, or a thing that seems good, takes \textbf{\ur{مجھے}} and \textbf{\ur{لگنا}}, and the verb follows the thing: \textbf{\ur{شہر}} is masculine, \textbf{\ur{اچھا} \ur{لگا}}. To ask:

\begin{quote}
\textbf{\ur{آپ} \ur{کو} \ur{فلم} \ur{کیسی} \ur{لگی؟}}
\end{quote}

\begin{quote}
\emph{āp ko film kaisī lagī?} — \textquotedblleft{}How did you like the film?\textquotedblright{}
\end{quote}

\subsection*{Your turn}
Say these aloud:

\begin{itemize}
\raggedright
  \item \textquotedblleft{}I was happy\textquotedblright{} — \emph{maiṅ khush thā}
  \item \textquotedblleft{}I was sad\textquotedblright{} — \emph{maiṅ udās thā}
  \item \textquotedblleft{}I was worried\textquotedblright{} — \emph{maiṅ pareshān thā}
\end{itemize}

\begin{tcolorbox}[breakable,colback=teal!4,colframe=teal!35!black,title={Before you move on}]
\textquotedblleft{}She was sad\textquotedblright{}? (\textbf{\ur{وہ} \ur{اداس} \ur{تھی}}.) \textquotedblleft{}I got scared\textquotedblright{}? (\textbf{\ur{مجھے} \ur{ڈر} \ur{لگا}}.) \textquotedblleft{}I liked the film\textquotedblright{}? (\textbf{\ur{مجھے} \ur{فلم} \ur{اچھی} \ur{لگی}}.)
\end{tcolorbox}
\section[\texorpdfstring{\ur{بننا} (bannā)}{بننا (bannā)}]{\ur{بننا} (bannā) — to become, to be made}

\label{lesson:UR-C272-banna}



\begin{quote}
Before the new one: say \textquotedblleft{}I'm happy\textquotedblright{}. (\emph{maiṅ khush hūṅ}.) Then \textquotedblleft{}to make\textquotedblright{}. (\emph{banānā}.) This lesson gives its partner.
\end{quote}

\subsection*{You'll want to know: \ur{بننا}}
\textbf{\ur{بننا}} — \emph{bannā} — \textquotedblleft{}to become; to be made\textquotedblright{}.

\begin{quote}
\textbf{\ur{وہ} \ur{استاد} \ur{بنی۔}}
\end{quote}

\begin{quote}
\emph{voh ustād banī.} — \textquotedblleft{}She became a teacher.\textquotedblright{}
\end{quote}

\begin{quote}
\textbf{\ur{یہ} \ur{گھر} \ur{پچھلے} \ur{سال} \ur{بنا۔}}
\end{quote}

\begin{quote}
\emph{yih ghar pichhle sāl banā.} — \textquotedblleft{}This house was built last year.\textquotedblright{}
\end{quote}

It has no object, so no \textbf{\ur{نے}}: \textbf{\ur{بنی}} follows her. Its partner \textbf{\ur{بنانا}} (\emph{banānā}), \textquotedblleft{}to make\textquotedblright{}, takes an object and \textbf{\ur{نے}}: \textbf{\ur{اس} \ur{نے} \ur{گھر} \ur{بنایا}} (\emph{us ne ghar banāyā}).

\begin{cousinweb}
\textbf{\ur{بننا}} (\emph{bannā}) is inherited. CDIAL traces it to Sanskrit \emph{vanati}, \textquotedblleft{}desires, gains, makes ready\textquotedblright{} (CDIAL 11260), and gives \textquotedblleft{}to be made\textquotedblright{} as a Hindi sense.
\end{cousinweb}

\subsection*{Your turn}
Say these aloud:

\begin{itemize}
\raggedright
  \item \emph{bannā}, then \emph{voh ustād banī}
  \item \textquotedblleft{}He became a teacher\textquotedblright{} — \emph{voh ustād banā}
\end{itemize}

\begin{tcolorbox}[breakable,colback=teal!4,colframe=teal!35!black,title={Before you move on}]
What does \textbf{\ur{بننا}} mean? (\textquotedblleft{}To become, to be made\textquotedblright{}.) Does it take \textbf{\ur{نے}}? (No: \textbf{\ur{بنانا}} does.)
\end{tcolorbox}
\section[\texorpdfstring{\ur{امید} (ummīd)}{امید (ummīd)}]{\ur{امید} (ummīd) — hope}

\label{lesson:UR-C272-ummid}



\begin{quote}
Before the new one: say \textquotedblleft{}I think she will come\textquotedblright{}. (\emph{mujhe lagtā hai ki voh āe gī}.) Now say it as a hope.
\end{quote}

\subsection*{You'll want to know: \ur{امید}}
\textbf{\ur{امید}} — \emph{ummīd} — \textquotedblleft{}hope\textquotedblright{}. It is feminine.

\begin{quote}
\textbf{\ur{مجھے} \ur{امید} \ur{ہے۔}}
\end{quote}

\begin{quote}
\emph{mujhe ummīd hai.} — \textquotedblleft{}I hope so.\textquotedblright{}
\end{quote}

\begin{quote}
\textbf{\ur{مجھے} \ur{امید} \ur{ہے} \ur{کہ} \ur{وہ} \ur{آئے} \ur{گی۔}}
\end{quote}

\begin{quote}
\emph{mujhe ummīd hai ki voh āe gī.} — \textquotedblleft{}I hope she'll come.\textquotedblright{}
\end{quote}

Word for word, \textquotedblleft{}to me there is hope\textquotedblright{}: \textbf{\ur{مجھے}}, then \textbf{\ur{کہ}} (\emph{ki}) and what you hope, in the future.

\begin{cousinweb}
\textbf{\ur{امید}} (\emph{ummīd}) is usually traced to Persian \emph{ummēd}, \textquotedblleft{}hope\textquotedblright{}.
\end{cousinweb}

\subsection*{Your turn}
Say these aloud:

\begin{itemize}
\raggedright
  \item \emph{ummīd}, then \emph{mujhe ummīd hai}
  \item \textquotedblleft{}I hope it won't rain\textquotedblright{} — \emph{mujhe ummīd hai ki bārish nahīṅ ho gī}
\end{itemize}

\begin{tcolorbox}[breakable,colback=teal!4,colframe=teal!35!black,title={Before you move on}]
What does \textbf{\ur{امید}} mean? (\textquotedblleft{}Hope\textquotedblright{}.) \textquotedblleft{}I hope she'll come\textquotedblright{}? (\textbf{\ur{مجھے} \ur{امید} \ur{ہے} \ur{کہ} \ur{وہ} \ur{آئے} \ur{گی}}.)
\end{tcolorbox}
\section[\texorpdfstring{\ur{بننا} \ur{چاہتی} \ur{ہوں} (bannā chāhtī hūṅ)}{بننا چاہتی ہوں (bannā chāhtī hūṅ)}]{\ur{بننا} \ur{چاہتی} \ur{ہوں} (bannā chāhtī hūṅ) — what I want to become}

\label{lesson:UR-C272-banna-chahti-hun}



\begin{quote}
Say \textquotedblleft{}she became a teacher\textquotedblright{}. (\emph{voh ustād banī}.) Then \textquotedblleft{}I want to go\textquotedblright{}. (\emph{maiṅ jānā chāhtā hūṅ}.) Put the two together.
\end{quote}

\subsection*{You'll want to know: I want to become}
\begin{quote}
\textbf{\ur{میں} \ur{استاد} \ur{ہوں۔}}
\end{quote}

\begin{quote}
\emph{maiṅ ustād hūṅ.} — \textquotedblleft{}I'm a teacher.\textquotedblright{}
\end{quote}

\begin{quote}
\textbf{\ur{میں} \ur{استاد} \ur{بننا} \ur{چاہتی} \ur{ہوں۔}}
\end{quote}

\begin{quote}
\emph{maiṅ ustād bannā chāhtī hūṅ.} — \textquotedblleft{}I want to become a teacher.\textquotedblright{}
\end{quote}

\textbf{\ur{چاہنا}} (\emph{chāhnā}) takes \textbf{\ur{بننا}} whole, as it takes \textbf{\ur{جانا}}: \textbf{\ur{بننا} \ur{چاہتی} \ur{ہوں}}, a woman speaking.

\begin{quote}
\textbf{\ur{وہ} \ur{نرس} \ur{بننا} \ur{چاہتا} \ur{ہے۔}}
\end{quote}

\begin{quote}
\emph{voh nars bannā chāhtā hai.} — \textquotedblleft{}He wants to become a nurse.\textquotedblright{}
\end{quote}

\subsection*{You'll want to know: my dream, my hope}
\begin{quote}
\textbf{\ur{میرا} \ur{خواب} \ur{استاد} \ur{بننا} \ur{ہے۔}}
\end{quote}

\begin{quote}
\emph{merā khvāb ustād bannā hai.} — \textquotedblleft{}My dream is to become a teacher.\textquotedblright{}
\end{quote}

\textbf{\ur{خواب}} (\emph{khvāb}) is masculine: \textbf{\ur{میرا}}. A hope looks ahead, so it takes the future:

\begin{quote}
\textbf{\ur{مجھے} \ur{امید} \ur{ہے} \ur{کہ} \ur{میں} \ur{استاد} \ur{بنوں} \ur{گی۔}}
\end{quote}

\begin{quote}
\emph{mujhe ummīd hai ki maiṅ ustād banūṅ gī.} — \textquotedblleft{}I hope I'll become a teacher.\textquotedblright{}
\end{quote}

\subsection*{Your turn}
Say these aloud:

\begin{itemize}
\raggedright
  \item \textquotedblleft{}I want to become a nurse\textquotedblright{} — \emph{maiṅ nars bannā chāhtī hūṅ}
  \item \textquotedblleft{}I want to become a player\textquotedblright{} — \emph{maiṅ khilāṛī bannā chāhtī hūṅ}
  \item \textquotedblleft{}I want to become an engineer\textquotedblright{} — \emph{maiṅ injīniyar bannā chāhtī hūṅ}
\end{itemize}

\begin{tcolorbox}[breakable,colback=teal!4,colframe=teal!35!black,title={Before you move on}]
\textquotedblleft{}I want to become a teacher\textquotedblright{}, a man speaking? (\textbf{\ur{میں} \ur{استاد} \ur{بننا} \ur{چاہتا} \ur{ہوں}}.) \textquotedblleft{}My dream\textquotedblright{}? (\textbf{\ur{میرا} \ur{خواب}}.) \textquotedblleft{}I hope she'll come\textquotedblright{}? (\textbf{\ur{مجھے} \ur{امید} \ur{ہے} \ur{کہ} \ur{وہ} \ur{آئے} \ur{گی}}.)
\end{tcolorbox}
\section[(review)]{Review — feelings and hopes}

\label{lesson:UR-R272-feelings-and-hopes}



\begin{quote}
Say \textquotedblleft{}a ticket, please\textquotedblright{} and \textquotedblleft{}my bag got lost\textquotedblright{}. (\emph{ek ṭikaṭ dījiye}, \emph{merā thailā kho gayā}.)
\end{quote}

\subsection*{Your turn}
\begin{itemize}
\raggedright
  \item \emph{Say it:} I'm happy today. Yesterday I was sad.
  \item \emph{Say it:} I'm afraid of dogs.
  \item \emph{Say it:} I liked this city. How did you like the film?
  \item \emph{Say it:} She became a teacher. I want to become a nurse.
  \item \emph{Say it:} My dream is to become a player. I hope I will become one.
  \item \emph{Write it:} \textbf{\ur{ڈر}} once, and point at \ur{ڈ} inside it
\end{itemize}

\begin{tcolorbox}[breakable,colback=teal!4,colframe=teal!35!black,title={Before you move on}]
\textquotedblleft{}I got scared\textquotedblright{}? (\textbf{\ur{مجھے} \ur{ڈر} \ur{لگا}}.) Now the other way round: in \textbf{\ur{مجھے} \ur{یہ} \ur{شہر} \ur{اچھا} \ur{لگا}}, which word sets \textbf{\ur{اچھا} \ur{لگا}}? (\textbf{\ur{شہر}}.) \textquotedblleft{}I hope\textquotedblright{}? (\textbf{\ur{مجھے} \ur{امید} \ur{ہے}}.)
\end{tcolorbox}
